\documentclass[11pt,a4paper,openany]{report}

% ============================================================
% PACKAGES
% ============================================================
\usepackage[a4paper,margin=1in]{geometry}
\usepackage{titlesec}
\usepackage{graphicx}
\usepackage{amsmath,amssymb}
\usepackage{booktabs}
\usepackage{longtable}
\usepackage{enumitem}
\usepackage{fancyhdr}
\usepackage{hyperref}
\usepackage{xcolor}
\usepackage{listings}
\usepackage{caption}
\usepackage{tocloft}
\usepackage{parskip}
\usepackage{array}

% ============================================================
% CODE LISTING STYLE
% ============================================================
\definecolor{codebg}{rgb}{0.96,0.96,0.96}
\definecolor{codekw}{rgb}{0.0,0.0,0.6}
\definecolor{codecom}{rgb}{0.35,0.55,0.35}
\definecolor{codestr}{rgb}{0.6,0.0,0.0}

\lstdefinestyle{pystyle}{
    backgroundcolor=\color{codebg},
    basicstyle=\ttfamily\footnotesize,
    keywordstyle=\color{codekw}\bfseries,
    commentstyle=\color{codecom}\itshape,
    stringstyle=\color{codestr},
    numberstyle=\tiny\color{gray},
    numbers=left,
    stepnumber=1,
    frame=single,
    rulecolor=\color{black!30},
    breaklines=true,
    breakatwhitespace=false,
    showstringspaces=false,
    tabsize=4,
    language=Python,
    captionpos=b
}

\lstset{style=pystyle}

% ============================================================
% PAGE STYLE
% ============================================================
\pagestyle{fancy}
\fancyhf{}
\fancyhead[L]{\small Solar Filament Segmentation 2026}
\fancyhead[R]{\small \leftmark}
\fancyfoot[C]{\thepage}
\renewcommand{\headrulewidth}{0.4pt}

% ============================================================
% TITLE FORMATTING — chapters flow inline (no forced page break)
% ============================================================
\titleclass{\chapter}{straight}
\titleformat{\chapter}[block]
  {\normalfont\Large\bfseries}
  {\chaptername\ \thechapter\quad}{0pt}{}
\titlespacing*{\chapter}{0pt}{24pt}{10pt}

% ============================================================
% HYPERREF SETUP
% ============================================================
\hypersetup{
    colorlinks=true,
    linkcolor=blue!60!black,
    urlcolor=blue!70!black,
    citecolor=green!40!black,
    pdftitle={Solar Filament Segmentation 2026 -- Technical Report},
    pdfauthor={Boobathi},
    pdfsubject={Instance Segmentation, YOLOv8, Solar Physics}
}

% ============================================================
% DOCUMENT
% ============================================================
\begin{document}

% ============================================================
% TITLE PAGE
% ============================================================
\begin{titlepage}
\centering
\vspace*{2cm}
{\Huge\bfseries Solar Filament Segmentation 2026\par}
\vspace{0.5cm}
{\Large An End-to-End Instance Segmentation Pipeline using YOLOv8l-Seg\\
on the MAGFiLO Dataset\par}
\vspace{1.5cm}
{\large A Technical Project Report\\
Based on the \texttt{solaris.py} Implementation\par}
\vspace{1.5cm}
{\large\itshape Prepared by: Boobathi\par}
\vspace{0.3cm}
{\large Department of Computer Science and Engineering\par}
\vspace{0.3cm}
{\large September 18, 2026\par}
\vspace{2cm}
\begin{abstract}
\noindent
Solar filaments are elongated, dark, thread-like structures observed in H-alpha images of the
Sun's chromosphere. Automated detection and segmentation of filaments is important for
space-weather monitoring, since filament eruptions are closely linked to Coronal Mass Ejections
(CMEs) that can affect satellites, power grids, and communication systems on Earth.

This report documents \textbf{Solaris}, an end-to-end deep learning pipeline implemented in
\texttt{solaris.py} that performs instance segmentation of solar filaments using the YOLOv8l-Seg
architecture, trained and evaluated on the MAGFiLO 1.0 dataset (Kaggle 2026 competition format).
The pipeline is designed to be portable across Kaggle, Google Colab, and local environments,
and covers five major stages:
\begin{enumerate}[label=(\arabic*)]
  \item Environment and path resolution across multiple hosting platforms.
  \item Dataset conversion from COCO-style polygon/RLE annotations into normalized YOLO
        segmentation labels.
  \item Model training of YOLOv8l-Seg at high resolution ($2048\times2048$) with domain-specific augmentation.
  \item Inference at high precision with strict confidence thresholding and duplicate suppression.
  \item Post-processing via priority-based overlap removal and conversion of binary masks into
        Run-Length Encoded (RLE) strings for Kaggle-style submission.
\end{enumerate}
Each function in \texttt{solaris.py} is explained line-by-line, with the underlying mathematics
(Dice score, IoU, mask normalization, RLE encoding) derived explicitly. The report closes with
an evaluation of the pipeline's design choices, its limitations, and directions for future improvement.
\end{abstract}
\end{titlepage}

% ============================================================
% TABLE OF CONTENTS
% ============================================================
\tableofcontents

% ============================================================
% CHAPTER 1 - INTRODUCTION
% ============================================================
\chapter{Introduction}

\section{What is a Solar Filament?}

A solar filament is a mass of relatively cool, dense plasma suspended above the Sun's surface by
magnetic field lines, appearing as a dark, thread-like structure when observed in H-alpha
(656.3\,nm) light against the brighter chromosphere. When the same structure is observed on
the solar limb (edge) rather than on the disk, it appears bright and is called a \textit{prominence}.
Filaments trace magnetic polarity inversion lines and are of major interest to space-weather
forecasters because their eruption is frequently associated with Coronal Mass Ejections (CMEs).

\section{Why Automate Filament Detection?}

Historically, filament identification has relied on manual annotation by solar physicists, which is
slow, subjective, and difficult to scale to the volume of daily full-disk H-alpha imagery produced
by observatories. An automated, pixel-accurate segmentation system enables:

\begin{itemize}
  \item Consistent, reproducible filament catalogs across years of data.
  \item Near real-time flagging of filaments for space-weather alerting.
  \item Quantitative morphological studies (length, area, orientation) at scale.
\end{itemize}

\section{Problem Statement}

Given a full-disk H-alpha solar image, the task is to produce a set of \textbf{pixel-accurate instance
segmentation masks}, one per filament, with \textbf{no overlapping pixels} between instances.
This is framed as a Kaggle-style competition (\texttt{filament-segmentation-2026}) using the
MAGFiLO 1.0 dataset, where predictions must be submitted as Run-Length Encoded masks in a
CSV file.

\section{Overview of the Solaris Pipeline}

The \texttt{solaris.py} script implements the complete pipeline in five logical stages:

\begin{center}
\renewcommand{\arraystretch}{1.3}
\begin{tabular}{>{\bfseries}c l l}
\toprule
Stage & Function(s) & Purpose \\
\midrule
1 & \texttt{resolve\_paths()} & Locate dataset and working directories \\
2 & \texttt{convert\_coco\_to\_yolo\_seg()} & Convert COCO annotations to YOLO format \\
3 & \texttt{train\_yolov8l\_seg()} & Train the segmentation model \\
4 & \texttt{run\_inference()} & Predict masks on test images \\
5 & \texttt{mask\_to\_rle()}, \texttt{remove\_instance\_overlaps()} & Post-process and encode submission \\
\bottomrule
\end{tabular}
\end{center}

% ============================================================
% CHAPTER 2
% ============================================================
\chapter{Environment and Dataset Resolution}

\section{The Portability Problem}

Machine learning notebooks are frequently developed on one platform (e.g.\ a local machine)
and then executed on another (Kaggle kernels, Google Colab). Each platform mounts datasets
at different filesystem paths. Rather than hard-coding a single path, Solaris implements an
auto-detection routine.

\section{The \texttt{resolve\_paths()} Function}

\begin{lstlisting}[caption={Environment auto-detection},label={lst:resolve}]
def resolve_paths():
    search_roots = [
        ".",
        "filament-segmentation-2026/MAGFiLO_1.0_Kaggle_2026",
        "/content",
        "/content/filament-segmentation-2026/MAGFiLO_1.0_Kaggle_2026",
        "/content/drive/MyDrive/kaggle_SF",
        "/kaggle/input/competitions/filament-segmentation-2026/...",
        "/kaggle/input/filament-segmentation-2026/..."
    ]
    for root in search_roots:
        if not os.path.exists(root):
            continue
        candidate_jsons = glob.glob(
            os.path.join(root, "**", "*Annotations*train*.json"),
            recursive=True)
        ...
\end{lstlisting}

\begin{enumerate}
  \item A list of candidate root directories is defined, covering local execution, Google Colab,
        Colab with Google Drive, and two Kaggle mount points.
  \item For each existing root, \texttt{glob.glob} with \texttt{recursive=True} searches for the COCO
        training JSON, a \texttt{train\_images} folder, and a \texttt{test\_images} folder.
  \item The first match for each category is kept so earlier entries take priority.
  \item Working directory is chosen by platform: \texttt{/kaggle/working} for Kaggle,
        \texttt{/content/working\_solaris} for Colab, or a local folder otherwise.
\end{enumerate}

\begin{lstlisting}[caption={Conditional submission path},label={lst:subpath}]
"output_sub": os.path.join(working_dir, "submission.csv")
             if "/kaggle/working" in working_dir
             else "submission.csv"
\end{lstlisting}

On Kaggle, only files written under \texttt{/kaggle/working} are captured as kernel outputs.

% ============================================================
% CHAPTER 3
% ============================================================
\chapter{The COCO Annotation Format}

\section{Why COCO?}

The MAGFiLO dataset ships annotations in the COCO JSON schema. A COCO file has three
top-level arrays:

\begin{center}
\renewcommand{\arraystretch}{1.3}
\begin{tabular}{l p{10cm}}
\toprule
\textbf{Key} & \textbf{Description} \\
\midrule
\texttt{images} & List of image records: \texttt{id}, \texttt{file\_name}, \texttt{width}, \texttt{height}. \\
\texttt{annotations} & Instance records with \texttt{image\_id}, \texttt{category\_id}, and \texttt{segmentation}. \\
\texttt{categories} & Class definitions (single class: \texttt{filament}). \\
\bottomrule
\end{tabular}
\end{center}

\section{Two Segmentation Encodings}

\begin{enumerate}
  \item \textbf{Polygon list} --- a list of flat coordinate arrays
        \texttt{[[x1, y1, x2, y2, \ldots, xn, yn]]}.
  \item \textbf{RLE dictionary} --- a dict with \texttt{counts} and \texttt{size} fields, decoded via
        \texttt{pycocotools.mask}.
\end{enumerate}

The conversion function branches on \texttt{isinstance(seg, list)} vs \texttt{isinstance(seg, dict)}.

% ============================================================
% CHAPTER 4
% ============================================================
\chapter{COCO $\to$ YOLO Label Conversion}

\section{Why Conversion is Necessary}

YOLOv8-Seg expects, for every image, a plain-text label file where each line is:
\[
  \langle\text{class\_id}\rangle\quad x_1\;y_1\;x_2\;y_2\;\cdots\;x_n\;y_n
\]
with coordinates normalized to $[0, 1]$:
\[
  x_\text{norm} = \frac{x_\text{pixel}}{w}, \qquad y_\text{norm} = \frac{y_\text{pixel}}{h}
\]

\section{The \texttt{convert\_coco\_to\_yolo\_seg()} Function}

\begin{lstlisting}[caption={Reproducible train/val split},label={lst:split}]
np.random.seed(seed)
image_ids = sorted(list(images.keys()))
np.random.shuffle(image_ids)
n_val     = int(len(image_ids) * val_ratio)
val_ids   = set(image_ids[:n_val])
train_ids = set(image_ids[n_val:])
\end{lstlisting}

\begin{lstlisting}[caption={Polygon normalization},label={lst:poly}]
if isinstance(seg, list):
    for poly in seg:
        if len(poly) < 6:
            continue
        coords      = np.array(poly).reshape(-1, 2)
        norm_coords = coords / [w, h]
        norm_coords = np.clip(norm_coords, 0.0, 1.0)
        flat        = " ".join([f"{x:.6f} {y:.6f}"
                                for x, y in norm_coords])
        label_lines.append(f"0 {flat}")
\end{lstlisting}

\begin{lstlisting}[caption={RLE decoding and contour extraction},label={lst:rle_dec}]
elif isinstance(seg, dict):
    m = mask_utils.decode(seg)
    contours, _ = cv2.findContours(
        m.astype(np.uint8),
        cv2.RETR_EXTERNAL,
        cv2.CHAIN_APPROX_SIMPLE)
    for cnt in contours:
        if len(cnt) < 3:
            continue
        coords = cnt.squeeze(1)
        ...
\end{lstlisting}

\begin{lstlisting}[caption={YOLO dataset descriptor (\texttt{data.yaml})},label={lst:yaml}]
yaml_content = f"""path: {os.path.abspath(output_dir)}
train: images/train
val:   images/val

names:
  0: filament
"""
\end{lstlisting}

% ============================================================
% CHAPTER 5
% ============================================================
\chapter{Image Preprocessing and Augmentation}

Solaris delegates resizing, normalization, and augmentation to the Ultralytics training loop:

\begin{itemize}
  \item \textbf{Mosaic augmentation} (\texttt{mosaic=0.5}): Four training images stitched into one
        composite with 50\% probability, exposing the model to varied scales.
  \item \textbf{Horizontal flip} (\texttt{fliplr=0.5}) and \textbf{Vertical flip} (\texttt{flipud=0.5}): each applied with 50\%
        probability. Full-disk solar images have no fixed ``up'' orientation, making flipping
        physically valid.
  \item \textbf{Resizing to \texttt{imgsz=2048}}: Images are letterboxed to $2048\times2048$,
        preserving the fine, thread-like structure of filaments that would be lost at YOLO's
        default 640-pixel resolution.
\end{itemize}

% ============================================================
% CHAPTER 6
% ============================================================
\chapter{YOLOv8l-Seg Architecture}

\section{Why an Instance Segmentation Model?}

Filaments can be numerous and closely spaced on a solar disk, so instance segmentation ---
which distinguishes individual object instances --- is required. YOLOv8-Seg extends YOLOv8
with a parallel mask-prediction branch.

\section{Architectural Components}

\begin{center}
\renewcommand{\arraystretch}{1.3}
\begin{tabular}{l p{9cm}}
\toprule
\textbf{Block} & \textbf{Role} \\
\midrule
Backbone (CSPDarknet) & Hierarchical feature extraction at decreasing spatial resolution. \\
Neck (PAN-FPN) & Bidirectional multi-scale feature fusion. \\
Detection Head & Anchor-free bounding box and class confidence prediction. \\
Mask Prototype Head & Predicts shared prototype masks. \\
Mask Coefficients & Per-instance weights that linearly combine the prototypes. \\
\bottomrule
\end{tabular}
\end{center}

The final instance mask:
\[
  M_i = \sigma\!\left(\sum_{k=1}^{K} c_{i,k} \cdot P_k\right)
\]

\section{Why the ``Large'' (l) Variant}

Filaments are thin, low-contrast structures that benefit from the deeper, wider backbone of
\texttt{yolov8l-seg.pt}. The training batch size is kept small (\texttt{batch=2}) to accommodate the
memory cost at $2048\times2048$ resolution.

% ============================================================
% CHAPTER 7
% ============================================================
\chapter{The Training Pipeline}

\section{The \texttt{train\_yolov8l\_seg()} Function}

\begin{lstlisting}[caption={Model training invocation},label={lst:train}]
model = YOLO("yolov8l-seg.pt")
results = model.train(
    data=yaml_path,
    epochs=epochs,
    imgsz=imgsz,
    batch=batch,
    device=0,
    workers=4,
    optimizer="AdamW",
    lr0=0.001,
    lrf=0.01,
    mosaic=0.5,
    fliplr=0.5,
    flipud=0.5,
    project=project,
    name="yolov8l_2048",
    save=True,
    verbose=True
)
\end{lstlisting}

\section{Parameter-by-Parameter Explanation}

\begin{center}
\renewcommand{\arraystretch}{1.3}
\begin{tabular}{l p{9cm}}
\toprule
\textbf{Parameter} & \textbf{Meaning} \\
\midrule
\texttt{data} & Path to the \texttt{data.yaml}. \\
\texttt{epochs} & Passes over the training set (default 30). \\
\texttt{imgsz=2048} & High resolution to preserve filament detail. \\
\texttt{batch=2} & Small batch due to memory cost of $2048\times2048$ on l-variant. \\
\texttt{device=0} & Train on GPU index 0. \\
\texttt{workers=4} & Four parallel CPU data-loading processes. \\
\texttt{optimizer="AdamW"} & Adam with decoupled weight decay. \\
\texttt{lr0=0.001} & Initial learning rate. \\
\texttt{lrf=0.01} & Final LR as fraction of \texttt{lr0} (cosine decay). \\
\texttt{mosaic, fliplr, flipud} & Augmentation probabilities (Chapter 5). \\
\texttt{project, name} & Output directory: \texttt{<project>/<name>/}. \\
\bottomrule
\end{tabular}
\end{center}

\section{The Loss Function}

\[
  \mathcal{L}_\text{total} = \lambda_\text{box}\mathcal{L}_\text{box}
                            + \lambda_\text{cls}\mathcal{L}_\text{cls}
                            + \lambda_\text{dfl}\mathcal{L}_\text{dfl}
                            + \lambda_\text{mask}\mathcal{L}_\text{mask}
\]

The mask component is per-pixel binary cross-entropy computed within the ground-truth
bounding box region.

\section{Checkpointing: \texttt{best.pt} vs \texttt{last.pt}}

Ultralytics saves \texttt{last.pt} (final epoch) and \texttt{best.pt} (best validation mAP epoch).
Solaris explicitly uses \texttt{best.pt}.

% ============================================================
% CHAPTER 8
% ============================================================
\chapter{Inference Pipeline}

\section{The \texttt{run\_inference()} Function}

\begin{lstlisting}[caption={Prediction call},label={lst:predict}]
result = model.predict(
    img_path,
    imgsz=2048,
    conf=0.30,
    iou=0.00,
    max_det=100,
    verbose=False
)[0]
\end{lstlisting}

\begin{center}
\renewcommand{\arraystretch}{1.3}
\begin{tabular}{l p{9cm}}
\toprule
\textbf{Parameter} & \textbf{Effect} \\
\midrule
\texttt{conf=0.30} & Suppress detections with confidence $<0.30$. \\
\texttt{iou=0.00} & Any box overlap triggers NMS suppression of the lower-confidence box. \\
\texttt{max\_det=100} & Upper bound on instances per image. \\
\bottomrule
\end{tabular}
\end{center}

\begin{lstlisting}[caption={Extracting raw prediction tensors},label={lst:tensors}]
result_masks_data = result.masks.data.cpu().numpy()
scores = result.boxes.conf.cpu().numpy() \
    if hasattr(result.boxes, 'conf') \
    else result.boxes.cls.cpu().numpy()
\end{lstlisting}

\texttt{result.masks.data} is a tensor of shape $(N, H_m, W_m)$ containing $N$ predicted binary
masks at the model's internal resolution.

% ============================================================
% CHAPTER 9
% ============================================================
\chapter{Priority-Based Overlap Removal}

\section{Why Overlaps Must Be Removed}

NMS operates on boxes, not pixel masks, so mask-level overlaps can persist. The competition
requires that no pixel belongs to more than one instance.

\section{The \texttt{remove\_instance\_overlaps()} Function}

\begin{lstlisting}[caption={Confidence-priority overlap elimination},label={lst:overlap}]
def remove_instance_overlaps(masks, scores):
    sorted_indices = np.argsort(scores)[::-1]
    sorted_masks   = masks[sorted_indices]

    accumulated_mask     = np.zeros(masks.shape[1:],
                                    dtype=np.uint8)
    non_overlapping_masks = []

    for mask in sorted_masks:
        clean_mask = mask & ~accumulated_mask
        non_overlapping_masks.append(clean_mask)
        accumulated_mask |= mask

    return np.array(non_overlapping_masks)
\end{lstlisting}

\begin{enumerate}
  \item Sort by confidence descending; the most confident detection gets first claim on pixels.
  \item \texttt{accumulated\_mask} tracks all pixels already claimed.
  \item \texttt{clean\_mask = mask \& \textasciitilde accumulated\_mask} removes pixels already claimed.
  \item \texttt{accumulated\_mask |= mask} adds the full footprint so later masks are clipped against it.
\end{enumerate}

\textbf{Worked example} (1D, 6 pixels):

\begin{center}
\renewcommand{\arraystretch}{1.3}
\begin{tabular}{lcccccc}
\toprule
& 0 & 1 & 2 & 3 & 4 & 5 \\
\midrule
Mask A (0.9) & 0 & 0 & 1 & 1 & 1 & 0 \\
Mask B (0.6) & 0 & 0 & 0 & 1 & 1 & 1 \\
Clean A      & 0 & 0 & 1 & 1 & 1 & 0 \\
Clean B      & 0 & 0 & 0 & 0 & 0 & 1 \\
\bottomrule
\end{tabular}
\end{center}

% ============================================================
% CHAPTER 10
% ============================================================
\chapter{Binary Mask to Run-Length Encoding}

\section{What is RLE?}

Run-Length Encoding compresses a binary mask by recording alternating run-lengths of
background and foreground pixels along a column-major (Fortran-order) traversal of the image.

\section{The \texttt{mask\_to\_rle()} Function}

\begin{lstlisting}[caption={RLE encoding via pycocotools},label={lst:rle_enc}]
def mask_to_rle(binary_mask):
    rle = mask_utils.encode(
        np.asfortranarray(binary_mask.astype(np.uint8)))
    rle["counts"] = rle["counts"].decode("utf-8")
    return rle["counts"]
\end{lstlisting}

\begin{enumerate}
  \item \texttt{np.asfortranarray} reorders memory to Fortran column-major order required by
        \texttt{pycocotools}.
  \item \texttt{mask\_utils.encode} walks the array and emits LEB128-style run lengths.
  \item \texttt{.decode("utf-8")} converts bytes to a plain string for the CSV cell.
\end{enumerate}

\section{Resizing Before Encoding}

\begin{lstlisting}[caption={Mask resize and trivial-fragment filtering},label={lst:resize}]
m_res = cv2.resize(m.astype(np.uint8), (w, h),
                   interpolation=cv2.INTER_NEAREST)
if m_res.sum() < 5:
    continue
\end{lstlisting}

\texttt{INTER\_NEAREST} is critical for binary masks; linear interpolation would introduce
non-binary intermediate values. Instances with fewer than 5 foreground pixels are discarded.

% ============================================================
% CHAPTER 11
% ============================================================
\chapter{Submission CSV Generation}

\section{Row Construction}

\begin{lstlisting}[caption={Building the submission DataFrame},label={lst:csv}]
submission_rows.append({
    "filament_id":      f"{file_stem}_{inst_idx}",
    "segmentation_rle": rle_str
})
inst_idx += 1
...
df_out = pd.DataFrame(submission_rows,
    columns=["filament_id", "segmentation_rle"])
df_out.to_csv(output_csv, index=False)
\end{lstlisting}

Each instance becomes one row identified by
\texttt{<image\_stem>\_<instance\_index>} (e.g.\ \texttt{img\_0451\_3}).
\texttt{index=False} matches Kaggle's expected schema.

\section{Sanity Reporting}

The script prints total instance count and number of unique images, giving a quick sanity
check before upload.

% ============================================================
% CHAPTER 12
% ============================================================
\chapter{Evaluation Metrics: Dice Score and IoU}

\[
  \text{Dice} = \frac{2|A \cap B|}{|A| + |B|}
              = \frac{2\,\text{TP}}{2\,\text{TP} + \text{FP} + \text{FN}}
\]
\[
  \text{IoU}  = \frac{|A \cap B|}{|A \cup B|}
              = \frac{\text{TP}}{\text{TP} + \text{FP} + \text{FN}}
\]
\[
  \text{Dice} = \frac{2\,\text{IoU}}{1 + \text{IoU}}
\]

where $A$ is the predicted mask, $B$ the ground-truth mask, TP the pixel overlap, FP
predicted-but-not-GT pixels, and FN GT-but-not-predicted pixels.

% ============================================================
% CHAPTER 13
% ============================================================
\chapter{The Main Entry Point and Orchestration}

\section{Command-Line Interface}

\begin{lstlisting}[caption={Argument parsing},label={lst:args}]
parser.add_argument("--train",   action="store_true")
parser.add_argument("--predict", action="store_true")
parser.add_argument("--model",   type=str,  default=None)
parser.add_argument("--epochs",  type=int,  default=30)
parser.add_argument("--imgsz",   type=int,  default=2048)
parser.add_argument("--batch",   type=int,  default=2)
\end{lstlisting}

Three invocation modes:
\begin{itemize}
  \item \textbf{Default}: convert + train, then infer using fresh weights.
  \item \textbf{\texttt{--train}}: convert + train only.
  \item \textbf{\texttt{--predict}}: skip training; use \texttt{--model} or fallback checkpoint search.
\end{itemize}

\section{Weight Fallback Search}

In predict-only mode, the script checks in order: the freshly trained checkpoint, a bundled
pretrained filename, and a Kaggle-dataset-mounted pretrained weights path.

% ============================================================
% CHAPTER 14 - RECOMMENDATIONS AND FUTURE WORK
% ============================================================
\chapter{Recommendations and Future Work}

This chapter presents concrete, actionable improvements to the Solaris pipeline that are
expected to push the Kaggle Dice score beyond the current 0.55 baseline. Each recommendation
includes a rationale and a ready-to-integrate code snippet.

% ----------------------------------------------------------
\section{Test-Time Augmentation (TTA)}

\textbf{Rationale:} Running the model on multiple augmented versions of the same image (horizontal
flip, vertical flip, and the original) and averaging the resulting masks typically improves recall
for thin, elongated structures like filaments.

\begin{lstlisting}[caption={TTA inference wrapper},label={lst:tta}]
import numpy as np
import cv2
from ultralytics import YOLO

def predict_with_tta(model_path, img_path,
                     imgsz=2048, conf=0.30):
    """
    Runs inference on original + 3 flipped versions,
    then merges masks by majority vote (>= 2/4 agree).
    """
    model = YOLO(model_path)
    img   = cv2.imread(img_path)
    h, w  = img.shape[:2]

    augmented_imgs = [
        img,                  # original
        cv2.flip(img,  1),   # horizontal flip
        cv2.flip(img,  0),   # vertical flip
        cv2.flip(img, -1),   # both axes
    ]
    flip_codes = [None, 1, 0, -1]

    vote_map = np.zeros((h, w), dtype=np.float32)

    for aug_img, flip_code in zip(augmented_imgs, flip_codes):
        tmp_path = "/tmp/tta_aug.jpg"
        cv2.imwrite(tmp_path, aug_img)

        result = model.predict(
            tmp_path, imgsz=imgsz,
            conf=conf, iou=0.00,
            max_det=100, verbose=False
        )[0]

        if result.masks is None:
            continue

        for raw_mask in result.masks.data.cpu().numpy():
            m = cv2.resize(raw_mask.astype(np.uint8),
                           (w, h),
                           interpolation=cv2.INTER_NEAREST)
            if flip_code is not None:
                m = cv2.flip(m, flip_code)   # undo flip
            vote_map += m.astype(np.float32)

    # Accept pixel if >= 2 out of 4 augments agree
    consensus_mask = (vote_map >= 2).astype(np.uint8)
    return consensus_mask
\end{lstlisting}

% ----------------------------------------------------------
\section{Multi-Scale Inference}

\textbf{Rationale:} Filaments span a wide range of apparent sizes. Running inference at two
resolutions (e.g.\ 1536 and 2048) and merging via mask union captures both fine-grained
details and broad structures.

\begin{lstlisting}[caption={Multi-scale inference and mask union},label={lst:multiscale}]
def predict_multiscale(model, img_path,
                       scales=(1536, 2048), conf=0.30):
    """Predict at multiple resolutions and union the masks."""
    img = cv2.imread(img_path, cv2.IMREAD_GRAYSCALE)
    h, w = img.shape
    combined = np.zeros((h, w), dtype=np.uint8)

    for scale in scales:
        result = model.predict(
            img_path, imgsz=scale,
            conf=conf, iou=0.00,
            max_det=100, verbose=False
        )[0]

        if result.masks is None:
            continue

        for raw_mask in result.masks.data.cpu().numpy():
            m = cv2.resize(raw_mask.astype(np.uint8),
                           (w, h),
                           interpolation=cv2.INTER_NEAREST)
            combined = np.maximum(combined, m)

    return combined
\end{lstlisting}

% ----------------------------------------------------------
\section{Confidence Threshold Tuning via Cross-Validation}

\textbf{Rationale:} The current \texttt{conf=0.30} was chosen empirically. A grid search over
the threshold on the validation split finds the value that maximises Dice without re-training.

\begin{lstlisting}[caption={Threshold sweep on validation set},label={lst:conf_sweep}]
import pandas as pd, json, numpy as np, cv2
from pathlib import Path
from pycocotools import mask as mask_utils

def sweep_conf_threshold(model_path, val_images_dir,
                         val_ann_json,
                         thresholds=(0.20, 0.25, 0.30,
                                     0.35, 0.40)):
    from ultralytics import YOLO
    model = YOLO(model_path)
    with open(val_ann_json) as f:
        gt_data = json.load(f)

    gt_by_img = {}
    for ann in gt_data["annotations"]:
        img_id = ann["image_id"]
        gt_by_img.setdefault(img_id, []).append(
            mask_utils.decode(ann["segmentation"])
            if isinstance(ann["segmentation"], dict)
            else None)

    records = []
    for conf in thresholds:
        dice_scores = []
        for img_info in gt_data["images"]:
            img_path = str(
                Path(val_images_dir) / img_info["file_name"])
            result = model.predict(
                img_path, imgsz=2048,
                conf=conf, iou=0.00,
                max_det=100, verbose=False)[0]

            h, w = img_info["height"], img_info["width"]
            pred_union = np.zeros((h, w), dtype=np.uint8)
            if result.masks is not None:
                for m in result.masks.data.cpu().numpy():
                    resized = cv2.resize(
                        m.astype(np.uint8), (w, h),
                        interpolation=cv2.INTER_NEAREST)
                    pred_union = np.maximum(pred_union, resized)

            gt_union = np.zeros((h, w), dtype=np.uint8)
            for gt_mask in gt_by_img.get(img_info["id"], []):
                if gt_mask is not None:
                    gt_union = np.maximum(gt_union, gt_mask)

            inter = (pred_union & gt_union).sum()
            dice  = (2 * inter /
                     (pred_union.sum() + gt_union.sum() + 1e-6))
            dice_scores.append(dice)

        records.append({"conf": conf,
                         "mean_dice": np.mean(dice_scores)})

    return pd.DataFrame(records).sort_values(
        "mean_dice", ascending=False)
\end{lstlisting}

% ----------------------------------------------------------
\section{Longer Training with Cosine LR Warm-Up}

\textbf{Rationale:} The current pipeline trains for 30 epochs with no warm-up. A 3-epoch linear
warm-up followed by cosine annealing over 60--100 epochs reduces early instability and allows
the model to find a better loss minimum.

\begin{lstlisting}[caption={Extended training with warm-up and cosine LR},label={lst:lr_warmup}]
def train_extended(yaml_path, epochs=80, imgsz=2048,
                   batch=2, project="magfilo_runs"):
    from ultralytics import YOLO
    model = YOLO("yolov8l-seg.pt")
    model.train(
        data=yaml_path,
        epochs=epochs,
        imgsz=imgsz,
        batch=batch,
        device=0,
        workers=4,
        optimizer="AdamW",
        lr0=0.001,
        lrf=0.005,        # tighter final LR (1/200 of lr0)
        warmup_epochs=3,  # linear warm-up
        cos_lr=True,      # cosine annealing schedule
        mosaic=0.5,
        fliplr=0.5,
        flipud=0.5,
        degrees=10.0,     # small random rotation
        project=project,
        name="yolov8l_2048_extended",
        save=True,
        verbose=True
    )
\end{lstlisting}

% ----------------------------------------------------------
\section{YOLOv8x-Seg (Extra-Large Variant)}

\textbf{Rationale:} The Extra-Large (x) variant has $\approx$68\,M parameters vs 44\,M for l.
If GPU memory allows (A100 40\,GB recommended at \texttt{batch=1}), switching to
\texttt{yolov8x-seg.pt} can improve mask mAP for thin, low-contrast objects.

\begin{lstlisting}[caption={Switching to YOLOv8x-Seg},label={lst:yolov8x}]
# Minimal change: replace model identifier
model = YOLO("yolov8x-seg.pt")   # was: yolov8l-seg.pt

model.train(
    data=yaml_path,
    epochs=50,
    imgsz=2048,
    batch=1,           # reduce batch for larger model
    device=0,
    optimizer="AdamW",
    lr0=0.0005,        # lower LR for larger model
    lrf=0.005,
    warmup_epochs=3,
    cos_lr=True,
    mosaic=0.5,
    fliplr=0.5,
    flipud=0.5,
    project="magfilo_runs",
    name="yolov8x_2048",
    save=True
)
\end{lstlisting}

% ----------------------------------------------------------
\section{Post-Processing: Morphological Closing}

\textbf{Rationale:} Small gaps in a predicted mask (from low-contrast mid-filament regions) can
split one true filament into multiple instances, reducing Dice. A $3\times3$ or $5\times5$
elliptical morphological closing bridges these gaps before RLE encoding.

\begin{lstlisting}[caption={Morphological closing to fill intra-filament gaps},label={lst:morph}]
def close_mask_gaps(binary_mask, kernel_size=5):
    """Apply morphological closing to fill small mask gaps."""
    kernel = cv2.getStructuringElement(
        cv2.MORPH_ELLIPSE,
        (kernel_size, kernel_size))
    closed = cv2.morphologyEx(
        binary_mask.astype(np.uint8),
        cv2.MORPH_CLOSE,
        kernel)
    return closed

# Insert before mask_to_rle() in run_inference():
m_res = cv2.resize(m.astype(np.uint8), (w, h),
                   interpolation=cv2.INTER_NEAREST)
m_res = close_mask_gaps(m_res, kernel_size=5)  # <-- NEW
if m_res.sum() < 5:
    continue
rle_str = mask_to_rle(m_res)
\end{lstlisting}




\end{document}
